\title{LongRoPE: Extending LLM Context Window Beyond 2 Million Tokens}

\begin{abstract}
A substantial context window is an advantageous characteristic for large language models (LLMs). Nevertheless, limitations such as the substantial expense of fine-tuning, the limited availability of lengthy textual data, and the emergence of problematic values with previously unseen token positions constrain present-day extended context windows to approximately 128k tokens.

This paper presents LongRoPE, a novel approach that, for the first time, extends the context window of pre-trained LLMs to an unprecedented \textbf{2048k} tokens. Remarkably, this is accomplished using no more than 1k fine-tuning steps, with training performed on sequences up to 256k tokens, all while preserving the original model’s performance for short contexts. The methodology relies on three principal innovations: (i) we discover and leverage two types of non-uniformities inherent in positional interpolation, utilizing an efficient search to establish superior fine-tuning initialization and enabling up to an 8$\times$ increase in context length even without additional fine-tuning; (ii) we propose a step-wise extension protocol—starting with fine-tuning an LLM at the 256k-token range, followed by a subsequent positional interpolation to further extend the fine-tuned model to a 2048k context window; (iii) we recalibrate LongRoPE on 8k-token sequences to restore its original short context capabilities. Comprehensive experiments using LLaMA2 and Mistral across multiple tasks underscore the effectiveness of our approach. Notably, models scaled up with LongRoPE retain their initial architectures, with only slight changes to the positional embedding layer, and are largely compatible with existing optimizations. The code will be released at \texttt{https://github.com/microsoft/LongRoPE}

\section{Introduction}

Although Large Language Models (LLMs) have achieved impressive performance across multiple tasks~\cite{gpt4,llama2}, they are constrained by a relatively small context window, such as the 4096-token boundary of LLaMA2~\cite{llama2}. When inputs exceed this window, the effectiveness of LLMs diminishes, as the model has not encountered such extended positions during its training. This limitation creates difficulties for critical applications, including in-context learning with a large number of samples~\cite{Huang2023FewerIM} and the operation of LLM-based agents~\cite{park2023generative,madaan2023selfrefine}.

Recent studies have demonstrated that it is possible to push the context window of a pre-trained LLM to approximately 128k tokens by applying fine-tuning on extended-length documents~\cite{longlora,pi,yarn,zhang2024soaring,liu2023scaling}. Nevertheless, scaling the context window even further faces three main challenges. First, assigning previously unused positional indices during extension causes numerous problematic outlier values, which result in out-of-distribution complications and subsequently impede the convergence of the fine-tuning process~\cite{pi}. This challenge is exacerbated when expanding from 4k to over 1000k positions, as this introduces more than 90\% novel indices. Second, successful fine-tuning demands the availability of sufficiently lengthy training texts. Yet, datasets containing very long sequences—particularly those surpassing 1000k tokens—are scarce. Furthermore, training with such ultra-long samples demands substantial computational resources, incurring very high costs in terms of training time and required GPUs. Third, with the adoption of extremely large context windows, the self-attention mechanism becomes diluted, since it has to distribute attention over a vast number of token positions. This diffusion diminishes performance on tasks involving shorter context lengths~\cite{pi}.

A common strategy for addressing the initial challenge involves applying interpolation to RoPE positional embeddings~\cite{rope,pi}, which effectively rescales novel position indices to align with those seen during pretraining, as depicted in Fig.\ref{fig:overview}. Position Interpolation (PI)~\cite{pi} implements a linear interpolation of the rotary angles in RoPE according to the desired extension ratio. Meanwhile, NTK~\cite{ntk,dynamicntk} proposes performing non-uniform (unequal) interpolation and extrapolation among the RoPE dimensions. In a different vein, YaRN~\cite{yarn} segments the RoPE dimensions into three classes determined by frequency, utilizing a combination of extrapolation, NTK, and linear interpolation for each group. Nevertheless, within Transformer models, positional embedding demonstrates a \emph{complex, uneven distribution of information entropy}. Current methodologies largely overlook such nuanced non-uniformity, which can cause information degradation and ultimately restricts the attainable context window size.

Section~\ref{sec:analysis} demonstrates through empirical evidence two principal observations: \textbf{(1)} Optimal positional interpolation should address both kinds of non-uniformity—namely, differences across RoPE dimensions and across token positions. Specifically, minimal interpolation is more suitable for lower RoPE dimensions and for tokens at earlier positions, although the ideal strategy depends on the desired sequence extension length. \textbf{(2)} Accounting for these non-uniformities in the interpolation process helps preserve crucial information inherent in the original RoPE, particularly for essential dimensions and early token positions. This approach reduces information loss due to positional interpolation and thereby yields improved initialization for subsequent fine-tuning. Additionally, this method enables sequence extension up to 8$\times$ even in settings without fine-tuning.

Building upon these observations, we propose \textbf{LongRoPE}, a robust approach capable of expanding the LLM context window to surpass 2 \emph{million} tokens. The design of LongRoPE incorporates three principal innovations. Firstly, {LongRoPE} capitalizes on positional interpolation by taking full advantage of the multidimensional and non-uniform structure inherent in positional encodings. This involves determining suitable rescale factors for the rotation angles in each RoPE dimension as a function of token positions. Given that the space of possible rescale factors grows exponentially with the desired extension ratio, {LongRoPE} employs an evolutionary search algorithm, enhanced by two optimization methods, to significantly improve the efficiency of the search. An illustration of the identified rescaled RoPE configuration is provided in Fig.~\ref{fig:overview}.

Subsequently, {LongRoPE} utilizes a resourceful, stepwise extension method to enable a context window of 2048k, thereby circumventing the necessity for direct fine-tuning on ultra-long text sequences, which are both rare and difficult to obtain. The approach initiates with a search for a suitable 256k length within the pre-trained LLM, followed by fine-tuning at this length. Owing to our non-uniform positional interpolation technique, which supports up to an 8$\times$ context length expansion without further fine-tuning, we perform an additional search to determine new RoPE rescale factors using the already fine-tuned, length-extended LLM. This procedure ultimately enables LLaMA2 and Mistral~\cite{mistral} models to support a context window size of 2048k.

In order to address potential declines in performance when processing shorter context windows, {LongRoPE} persistently refines the RoPE rescaling parameters within the extended LLM. In analogy to the expansion from a 256k to a 2048k window, our approach involves reducing the context window to 4k and 8k on the LLM that has been fine-tuned at 256k, leveraging our search algorithm to minimize positional interpolation. For inference, whenever the input sequence is shorter than 8k, we replace the RoPE parameters with the optimized rescale factors identified through our search.

Comprehensive evaluations involving multiple LLMs and a range of long-context tasks confirm the efficacy of our approach. Our results indicate that LongRoPE sustains low perplexity levels over evaluation lengths from 4k up to 2048k, attains passkey retrieval accuracy exceeding 90\%, and matches baseline performance on established benchmarks built for the 4096 context window. The LongRoPE framework is compatible with any LLM architecture utilizing RoPE embeddings. We intend to make both our codebase and the LongRoPE-2048k models publicly available.

\section{Non-uniformity in Positional Interpolation}
\label{sec:analysis}

\subsection{Preliminary}

Transformer architectures depend on explicit positional encodings—typically provided via position embeddings—to encode the sequence order of input tokens. In this study, we examine the RoPE~\cite{rope} positional embedding method, a prevalent choice in contemporary LLMs. For a token situated at the positional index $n$, its RoPE representation can be concisely expressed as:

\begin{equation}
		\fontsize{7.8}{7.8} \selectfont
	\label{eq:rope}
	\left[ cos(n\theta_0), sin(n\theta_0), cos(n\theta_1), \cdot\cdot\cdot, cos(n\theta_{d/2-1}), sin(n\theta_{d/2-1}) \right]   
\end{equation}

Here, $d$ stands for the embedding dimensionality, and $n\theta_i$ denotes the rotary angle assigned to the token located at position $n$. The term $\theta_i=\theta^{-2i/d}$ indicates the corresponding rotation frequencies. By default, RoPE uses $\theta=10000$ as the base for these frequencies.

\textbf{\emph{Context window extension ratio $s$} and positional interpolation}. We define $s$ as  the ratio of extended context length $L'$ to the original length $L$: $s=\frac{L'}{L}$. 

To increase the context window from $L$ to $L'$, existing positional interpolation strategies typically recommend reducing the rotation frequencies $\theta_i$ using the extension ratio $s$. Introducing $\beta=\theta^{2/d}$ and denoting the actual scaling factor corresponding to $s$ as $\lambda$, we can present a general formulation encompassing these positional interpolation methods:

\begin{equation}	
	\fontsize{7.5}{7.5} \selectfont
	\label{eq:unified_rope}
	\begin{aligned}
		\left[cos\left(\frac{n}{\lambda(\beta)^0}\right), sin \left(\frac{n}{\lambda(\beta)^0}\right), cos\left(\frac{n}{\lambda(\beta)^1}\right), \ldots,  sin \left(\frac{n}{\lambda(\beta)^{d/2-1}}\right) \right]   
	\end{aligned}
\end{equation}

\textbf{Linear positional interpolation (PI)}. PI~\cite{pi} employs a linear rescaling of position indices when staying within the model's original training sequence length. To achieve an extension ratio of $s$, the rotation angles associated with each position are uniformly scaled by a factor of $\lambda=s$ over all RoPE dimensions. This strategy compresses positional signals, making the representations of adjacent tokens more similar and thus diminishing the model's capacity to differentiate between nearby positions. Consequently, PI is often less effective at larger extension ratios.

\textbf{NTK-based interpolation and extrapolation}. The works of~\cite{ntk,dynamicntk} approach RoPE by viewing it through the lens of information encoding, utilizing the Neural Tangent Kernel (NTK) framework~\cite{jacot2018neural,tancik2020fourier}. To address the challenge of congested positional indices encountered in PI, they propose a method that redistributes the demands of interpolation across the various dimensions of RoPE. Specifically, their approach involves applying less scaling to lower (higher frequency) dimensions and greater scaling to higher (lower frequency) dimensions. This technique facilitates both the interpolation and extrapolation of positional encodings, following the transformation $\lambda=s^{i}$. The refined dynamic NTK method~\cite{dynamicntk} further adapts the extension factor dynamically at each position as a function of the input sequence length. In contrast to PI, which relies on fine-tuning to operate effectively, NTK-based strategies allow for the extension of context windows without the need for additional training, though typically only up to a 4$\times$ extension ratio.

\textbf{YaRN}~\cite{yarn} presents a notable advancement in the effectiveness of positional interpolation. Its approach separates the RoPE dimensions into three categories based on frequency, assigning a unique interpolation method to each group. Specifically, extrapolation ($\lambda$=1) is applied to high frequency dimensions, linear interpolation (PI) is used for low frequency dimensions, and the NTK technique is reserved for the intermediate frequencies. The central innovation of YaRN is this frequency-based partitioning of RoPE dimensions; however, this process currently relies on manual tuning and empirical observation by researchers, which could hinder optimal results when deploying to newly developed LLMs.

\subsection{Study on Non-uniform Positional Interpolation}

Building on ideas from NTK and YaRN, we observe that their improvements stem from leveraging non-linearity—particularly through treating RoPE dimensions with frequency-specific interpolation and extrapolation strategies. Nonetheless, prevailing non-linear methods predominantly depend on handcrafted heuristics. This observation prompts two key inquiries: (1) Does existing positional interpolation represent the best approach? (2) Might there be additional non-linearities that have not yet been investigated?

In addressing these questions, we employ evolution search (refer to Sec.~\ref{sec:method}) to identify improved non-uniform positional interpolation strategies for LLaMA2-7B. This search process is steered by perplexity metrics, evaluated over five randomly selected samples from the PG19~\cite{pg19} validation set. Our experimental results yield the following principal insights.

\textbf{Finding 1}: \textit{RoPE dimensions exhibit substantial non-uniformities, which are not effectively handled by current positional interpolation methods.}

We perform a search for the optimal $\lambda$ value for each RoPE dimension as described in Eq.~\ref{eq:unified_rope}. Table~\ref{tbl:exp1} reports the perplexity results of LLaMA2-7B across various approaches on the PG19 and Proof-pile~\cite{proof-pile} test datasets, all evaluated without any fine-tuning. Our optimized approach yields notable gains, indicating that both linear (PI) and non-uniform (Dynamic-NTK, YaRN) interpolation strategies are not optimal. In particular, YaRN lags behind PI and NTK on PG19, as it fails to span the intended context window for LLMs that have not undergone fine-tuning. For instance, with an 8k context size, YaRN’s perplexity rises steeply after reaching 7k.

In our investigation, the rescaled coefficients $\lambda$ in Eq.~\ref{eq:unified_rope} exhibit variability, in contrast to the constant scale $s$ applied in PI, the formula derived from NTK, and the group-wise computation used by YaRN. This introduction of non-uniform scaling leads to notable enhancements in the language modeling capabilities of LLaMA2, specifically reflected in reduced perplexity scores for 8k and 16k context lengths, even without additional fine-tuning. This improvement arises because the modified positional embeddings more faithfully reconstruct the original RoPE structure, particularly in critical dimensions, thereby easing the model's challenge of distinguishing nearby token positions.

\textbf{Finding 2}: \textit{RoPE for the initial tokens in the input sequence should be extrapolated with less interpolation.} 

We propose that for the initial $\hat{n}$ tokens within the input sequence, the RoPE mechanism should apply reduced interpolation. This stems from the fact that these early tokens are assigned high attention scores, which, as reported by Streaming LLM~\cite{streamingllm} and LM-Infinite~\cite{han2023lminfinite}, makes them particularly influential within attention modules. To test this proposition, we experimented with expanding the context length to 8k and 16k using both PI and NTK approaches, selectively disabling interpolation for the first $\hat n$ tokens ($\hat n = 0, 2, \ldots, 256$). Setting $\hat n = 0$ corresponds to the standard PI and NTK configurations. Results, summarized in Table~\ref{tbl:tokenposition}, yield two primary insights: \textbf{(1)} maintaining the original representations for the earliest tokens (i.e., skipping position interpolation) leads to measurable performance gains; \textbf{(2)} the ideal choice of $\hat n$ varies depending on the desired extension length of the context window.

\textbf{Finding 3}: \textit{Non-uniform positional interpolation effectively extends LLM context window in both fine-tuning and non-fine-tuning settings.} 

Although our results demonstrate that the proposed non-uniform position interpolation notably enhances extension performance at 8k and 16k contexts without fine-tuning, fine-tuning remains necessary for even longer context lengths. Therefore, we conduct fine-tuning of LLaMA2-7B using our discovered RoPE configuration for a 64k context window (refer to the Appendix for experimental details). As indicated in Table~\ref{tbl:finetune}, our approach achieves substantially superior results compared to PI and YaRN, both prior to and following the fine-tuning of LLaMA2-7B. This superiority is attributed to the effective exploitation of non-uniform positional interpolation, which mitigates information loss and establishes a more advantageous starting point for subsequent fine-tuning.

\textbf{Summary}. This work identifies two forms of non-uniformity: heterogeneity in RoPE dimensions and token locations. Leveraging these aspects within positional interpolation considerably enhances LLM performance on extended context windows.

\section{LongRoPE}

\label{sec:method}
Building upon these observations, we propose LongRoPE, which initially incorporates a novel and efficient search strategy to leverage both types of non-uniformity. Subsequently, this approach is utilized to expand the context window of LLMs to support over 2 million tokens.

\subsection{Problem Formulation}

The presence of these two non-uniformities significantly enlarges the solution space and complicates the optimization process. To manage this, we conceptualize the challenge of multidimensional non-uniform position interpolation optimization as a search problem.

Considering an LLM intended for a context window of size $L'$, and a set of long input documents $\mathbf{X}$ such that each document $\mathbf{x}\in\mathbf{X}$ contains more tokens than $L'$, we represent the initial rotary angle for the $i^{th}$ dimension in the RoPE embedding at token index $n$ as $\frac{n}{\beta^i}$. The corresponding optimization task is expressed as:

\begin{equation}
\begin{gathered}
		\label{eq:problem}

\underset{\mathbf{x}\in\mathbf{X};\,|\mathbf{x}|\ge L'}{\operatorname{arg\,min}}\,\mathcal{L}\left(\text{LLM}(\text{RoPE}, \mathbf{X})\right),\quad \text{with}
\\
\scalebox{0.93}{
$\underset{\begin{subarray}{l} i=0,\ldots,\frac{d}{2}-1; \\ n\in[0,|\mathbf{x}|); \end{subarray}}{\text{RoPE}_{(n)}} = \left[\, \ldots, \cos\left(\mathbb{I}(\hat{\lambda}_{i},\hat{n})\, \frac{n}{\beta^i}\right),\; \sin\left(\mathbb{I}(\hat{\lambda}_{i},\hat{n})\,\frac{n}{\beta^i}\right),\ldots \,\right] $}
\\
\text{where } \mathbb{I}(\hat{\lambda}_{i},\hat{n})= 
\begin{cases}
    1 & \text{if } n < \hat{n} \\
    {\frac{1}{\lambda}_{i}} & \text{if } n \geq \hat{n}
\end{cases}
\end{gathered}
\end{equation}

In this context, we define a collection of rescaling factors, $\mathbb{I}(\hat{\lambda}_{i},\hat{n})$, designed to address both identified types of non-uniformity. Here, $\hat{\lambda}_i$ represents the non-uniformity corresponding to specific RoPE dimensions, while $\hat{n}$ refers to non-uniformities in token positions. The function $\mathbb{I}(\hat{\lambda}_{i},\hat{n})$ adjusts the rotational angle for the $i^{th}$ dimension of RoPE by employing $\hat{\lambda}_i$ as a scaling factor, with $\hat{n}$ specifying the threshold for token positions. For token positions less than $\hat{n}$, i.e., the initial $\hat{n}-1$ tokens, the original rotary angle $\frac{n}{\beta^i}$ is retained without modification by the scaling factor. When token positions satisfy $n\ge\hat{n}$, the rescaling factor becomes operative.

Suppose we are given a desired context window size of $L'$. Our task is to determine the most suitable set of rescale factors ($\mathbb{I}(\hat{\lambda}_{0},\hat{n})$, $\mathbb{I}(\hat{\lambda}_{1},\hat{n})$, ...$\mathbb{I}(\hat{\lambda}_{i},\hat{n})$, etc.) across all RoPE dimensions, indexed from $1$ through $d$. By appropriately rescaling the $\text{RoPE}$ in the target $\text{LLM}$, the model aims to minimize the next token prediction loss, denoted as $\mathcal{L}$ (also understood as perplexity), when processing input sequences $\mathbf{X}$ of length $L'$.

\subsection{Searching the Non-uniform Position Interpolation}

To address the challenge posed by Eq.~\ref{eq:problem}, we propose a straightforward yet powerful approach that identifies optimal RoPE rescale factors, thereby maximizing the potential of multidimensional non-uniformities present within position embeddings.

\textbf{Search space}. We construct an expansive search space that captures both types of non-uniformities. An overview of this search space is presented in Table~\ref{tbl:searchspace}. More precisely, our framework enables the assignment of individualized rescale factors to each RoPE embedding dimension. To simplify the architecture of the search space, the search is performed over $\lambda_i$ and $\hat{n}$, as opposed to directly searching for $\mathbb{I}(\hat{\lambda}_{i},\hat{n})$, with $\hat{\lambda}_{i}=1/\lambda_i$. Table~\ref{tbl:searchspace} details that $\lambda_i$ ranges from a baseline value of 1.0, representing pure extrapolation, up to a maximum bound of $s\times1.25$, which supports interpolation scenarios exceeding the PI configuration, with increments set at 0.01. Here, $s$ denotes the proportion by which the context window is intended to be expanded.

The parameter $\hat{n}$ determines how many leading token positions preserve their original positional encoding by forgoing interpolation and directly employing the RoPE embedding. In practice, we search over the following set of possible values for $\hat{n}$: \{0, 1, 2, 4, 8, 12, 16, 20, 24, 28, 32, 64, 128, 256\}. Setting $\hat{n}=0$ implies that every token position is assigned a rescale factor found during the search process.

\textbf{Evolution-based search}. The search space outlined in Table~\ref{tbl:searchspace} encompasses a wide array of positional interpolation configurations, making efficient navigation through it highly complex. For instance, setting $s=4\times$ for extension yields $400^{128/2}\times14 = 4\times10^{167}$ possible combinations. As the extension ratio increases, this search space grows at an exponential rate. To manage this complexity, we employ evolution search~\cite{guo2020single} and incorporate two optimization strategies to markedly enhance search efficiency. Algorithm~\ref{alg:search} presents an overview of the complete search process.

\textit{Optimized initial population generation}. Rather than generating all $P$ rescale factors at random for the initial population, we explicitly include the three RoPE rescale factors derived from PI, NTK, and YaRN as individual members. The remaining $P$-3 individuals are produced by applying random mutations, with probability $p$, to these three rescale factors.

\textit{Monotonically non-decreasing constraint}. Once the initial population is produced, LLM perplexity is measured for every individual. This involves applying each candidate’s specific RoPE rescale parameters to the target LLM and evaluating the perplexity over input $\mathbf{X}$. The best $k$ individuals, ranked by their perplexity scores, are chosen as parents for the subsequent evolutionary steps. Nevertheless, due to the immense size of the search space, unrefined mutation and crossover operations often lead to inferior candidates and, consequently, a surplus of unnecessary perplexity evaluations. This inefficiency becomes especially pronounced when $L'$ is substantial, as each perplexity computation is computationally intensive.

To tackle this issue, we enforce a constraint ensuring that the sampled RoPE rescaling factors exhibit non-decreasing monotonicity: $\lambda_i \le \lambda_{i+1}$. Only those RoPE configurations that comply with this requirement are used for perplexity evaluation within the LLM, thereby substantially curtailing the search space. Concretely, we demand that $\lambda_i$ should monotonically increase alongside the RoPE dimension (where $i=0, \ldots, 63$). This constraint on dimensional monotonicity is grounded in NTK theory~\cite{jacot2018neural,tancik2020fourier,ntk}, which indicates that lower-dimensional, higher-frequency components necessitate less interpolation (corresponding to smaller $\lambda_i$ values), whereas higher-dimensional, lower-frequency components allow for more interpolation (implying larger $\lambda_i$).

\begin{algorithm}[t]
	\small
	\caption{The search algorithm}
	\textbf{Input:} target LLM, input samples $\mathbf{X}$, population size $P$, mutation size $N_1$, crossover size $N_2$, maximum iterations $\mathcal{T}$, mutation probability $p$\\

	\begin{algorithmic}[1]
		\label{alg:search}
		\STATE Initialize Top-k as an empty set;
		\STATE Set $\text{P}_0$ = \textit{Initialize\_population\_with\_optimization}($P$, $\mathbf{X}$, $p$);
		\FOR{$i = 1$ to $\mathcal{T}$}
		\STATE Execute \textit{Compute\_perplexity}(LLM, $\text{P}_{i-1}$, $\mathbf{X}$);
		\STATE Top-k $\leftarrow$ \textit{Update\_Topk}(Top-k, $\text{P}_{i-1}$);
		\STATE $\text{P}_{mutation}$ = \textit{Mutation\_with\_mono\_constraint}(Top-k, $N_1$, $p$);
		\STATE $\text{P}_{crossover}$ = \textit{Crossover\_with\_mono\_constraint}(Top-k, $N_2$);
		\STATE $\text{P}_i$ = $\text{P}_{mutation} \cup \text{P}_{crossover} \cup$ Top-k;
		\ENDFOR
		\STATE Output the Top-k member with the minimal perplexity;
	\end{algorithmic}
\end{algorithm}

\textbf{8$\times$ extension without fine-tuning}. Through the use of evolutionary search, our approach successfully discovers non-uniform RoPE rescale factors that retain critical dimensions and positional encodings, thereby reducing the informational degradation caused by interpolation. As shown in Fig.\ref{fig:non-ft}, this strategy enables us to expand LLaMA2’s context length from 4k to 32k tokens without necessitating fine-tuning. By comparison, prior techniques like PI and non-uniform NTK and YaRN exhibit marked increases in perplexity beyond a 2$\times$ extension.

\subsection{Extending LLM Context Window to 2048K}
\label{sec:extendto2048k}

\textbf{Progressive extension to 2048k}. In this section, we present a strategy to expand the context window of existing pre-trained LLMs from the standard 4k up to beyond 2048k tokens. Our results indicate that non-uniform positional interpolation enables an 8$\times$ window increase without necessitating additional fine-tuning. For even greater extensions (for example, a 512$\times$ expansion), fine-tuning becomes unavoidable. One approach involves identifying suitable RoPE scaling factors at the target 2048k context length, followed by fine-tuning the model. Nonetheless, this approach is hindered by the substantial computational costs associated with training. Furthermore, in our experiments, effective fine-tuning at such high extension ratios proves to be quite challenging (refer to Appendix for details).

Fortunately, LongRoPE proves to be successful when applied to both baseline and fine-tuned extended LLMs. Consequently, we propose a streamlined and incremental approach that enables the model to reach the desired 2048k target by performing only 1k fine-tuning iterations, all conducted within a training length of 256k.

$\Diamond$ \textit{Expanding pre-trained LLMs to a 256k context window using LongRoPE search}. For illustration, with LLaMA2, we perform search operations aimed at achieving context window sizes of 128k and 256k. At these scales, the context length is increased by a factor of 32$\times$ and 64$\times$ correspondingly.

$\Diamond$ \textit{Fine-tuning to 256k}. Subsequently, we adapt the pre-trained LLM to reach a context window of 256k tokens. To do so, we begin by fine-tuning LLaMA2 for 400 steps with the RoPE rescaled factors set for 128k. After this initial phase, we switch the RoPE rescaled factors to those for 256k on the obtained checkpoint and continue with a further 600 fine-tuning steps. This staged approach is more efficient than commencing the fine-tuning process directly at 256k.

$\Diamond$ \textit{LongRoPE Search for Scaling Extended LLMs to 2048k}. In the concluding stage, we conduct an additional search process utilizing the 256k-length LLM that has already been fine-tuned. This approach allows the model’s context window to be dramatically enlarged to 2048k, achieving this significant expansion without any extra fine-tuning procedures. The total extension achieved corresponds to a $512\times$ increase.

\textbf{Recovery within reduced context windows}. Upon expanding the context window to a very large 2048k length, there is a notable decrease in model performance for the initial context window size. This phenomenon has been documented with positional interpolation~\cite{pi}, since embedding positions in higher-dimensional space results in those corresponding to the original context window occupying a considerably confined area, thereby harming language model efficacy. When applying an extension factor of 512$\times$, the original 4k context window's positional representations become especially congested.

To address this issue, we conduct an additional evolution search on the expanded LLM, specifically tuning the RoPE rescale factors for relatively short context lengths such as 4k and 8k. Since positional interpolation demands are lower for these shorter sequences, we constrain the search space by reducing the upper bound for $\lambda$. At inference time, the LLM adapts the relevant RoPE rescale factors in real time.

\section{LongRoPE}

\label{sec:method}
Building upon these insights, we propose LongRoPE—a method that initially incorporates an effective search algorithm designed to take full advantage of the identified non-uniformities, and subsequently leverages this algorithm to expand the context window of LLMs to lengths exceeding 2 million tokens.

\subsection{Problem Formulation}

These two forms of non-uniformity significantly expand the range of possible solutions and add layers of difficulty to the optimization process. To tackle this challenge, we recast the multidimensional non-uniform position interpolation optimization as a search-based problem.

Given an LLM designed for a context window of size $L'$ and a collection of long documents $\mathbf{X}$—with each document $\mathbf{x}\in\mathbf{X}$ containing more tokens than $L'$—let us represent the initial rotary angle assigned to the $i^{th}$ RoPE embedding dimension at token position $n$ by $\frac{n}{\beta^i}$. Under this setup, the corresponding optimization problem can be stated as:

\begin{equation}
\begin{gathered}
		\label{eq:problem}

\underset {\mathbf{x}\in\mathbf{X}; \, |\mathbf{x}|\ge L'}{ \operatorname {arg\,min} } \,  \mathcal{L}\left( \text{LLM} (\text{RoPE}, \mathbf{X})\right), \quad \text{with}
\\
\scalebox{0.93}{
$\underset {\begin{subarray}{l} i=0,\dots,\frac{d}{2}-1;\\ n\in[ 0,|\mathbf{x}|); \end{subarray}}{\text{RoPE}_{(n)}} = {\left[\dots,\cos\left(\mathbb{I}(\hat{\lambda}_{i}, \hat{n}) \frac{n}{\beta^i}\right),\; \sin\left(\mathbb{I}(\hat{\lambda}_{i}, \hat{n}) \frac{n}{\beta^i}\right),\dots \right]}$}\\
\text{where} \; \mathbb{I}(\hat{\lambda}_{i}, \hat{n}) = \begin{cases}
1 &\text{\quad if $n<\hat{n}$}\\
\frac{1}{\lambda_{i}} &\text{\quad if $n\geq\hat{n}$}
\end{cases}
	\end{gathered}
\end{equation}

In this approach, we define a group of rescale factors, denoted as $\mathbb{I}(\hat{\lambda}_{i},\hat{n})$, to address two distinct types of non-uniformity. Here, $\hat{\lambda_i}$ represents the variation present in RoPE dimensions, while $\hat{n}$ indicates the token position threshold. The function $\mathbb{I}(\hat{\lambda}_{i},\hat{n})$ specifically adjusts the rotation angle associated with the $i^{th}$ RoPE dimension, with $\hat{\lambda}_i$ serving as the scaling parameter and $\hat{n}$ marking the boundary in token positions. For token positions less than $\hat{n}$, meaning the first $\hat{n}-1$ positions, the rotation angles remain unchanged, retaining the original RoPE form $\frac{n}{\beta^i}$. When token positions satisfy $n \ge \hat{n}$, the rescale factor is activated and applied to the rotation angle.

For a specified context window length $L'$, the goal is to determine the set of optimal rescale factors ($\mathbb{I}(\hat{\lambda}_{0},\hat{n})$, $\mathbb{I}(\hat{\lambda}_{1},\hat{n})$, ... $\mathbb{I}(\hat{\lambda}_{i},\hat{n})$...) corresponding to each of the $1^{st}$ through $d^{th}$ RoPE dimensions. Implementing these rescaling adjustments within the target $\text{LLM}$’s $\text{RoPE}$ mechanism should enable the model to minimize its next token prediction loss, $\mathcal{L}$ (as measured by perplexity), on input sequences $\mathbf{X}$ of token length $L'$.

\subsection{Searching the Non-uniform Position Interpolation}

To address the issue outlined in Eq.~\ref{eq:problem}, we propose an efficient and straightforward approach that identifies optimal RoPE rescaling factors. This approach is designed to take full advantage of the complex multidimensional variations present in position embeddings.

\textbf{Search space}. Our method features an expansive search space capable of capturing both types of non-uniformities. Table~\ref{tbl:searchspace} summarizes this search space configuration. Rather than assigning a universal scaling factor, we assign a unique rescale factor to each dimension within the RoPE embedding. To facilitate a more tractable search, our parameterization uses $\lambda_i$ and $\hat{n}$ instead of directly optimizing over $\mathbb{I}(\hat{\lambda}_{i},\hat{n})$, given that $\hat{\lambda}_{i}=1/\lambda_i$. As detailed in Table~\ref{tbl:searchspace}, the search for $\lambda_i$ is constrained to the interval starting from 1.0 (representing simple extrapolation) up to $s\times1.25$ (allowing for a more expansive interpolation than what is used in PI), using increments of 0.01. Here, $s$ denotes the factor by which the desired context window is to be extended.

The parameter $\hat{n}$ determines how many of the earliest token positions keep their original RoPE embeddings instead of undergoing position interpolation. In practice, $\hat{n}$ is chosen from the set \{0, 1, 2, 4, 8, 12, 16, 20, 24, 28, 32, 64, 128, 256\}. Setting $\hat{n}=0$ means every token position applies the optimized rescale factors.

\textbf{Evolution-based search}. The search space presented in Table~\ref{tbl:searchspace} encompasses a vast array of positional interpolation configurations, making exhaustive search impractical. For instance, increasing the extension factor to $s=4\times$ yields $400^{128/2}\times14=4\times10^{167}$ possible options. As the extension ratio grows, the size of the search space escalates exponentially. To tackle this issue, we adopt evolution search~\cite{guo2020single} and incorporate two optimization strategies that substantially improve the efficiency of the search. The entire search process is outlined in Algorithm~\ref{alg:search}.

\textit{Optimized initial population generation}. Rather than relying solely on random initialization for the population of $P$ rescale factors, we explicitly insert the RoPE rescale factors used by PI, NTK, and YaRN as specific individuals in the initial set. The rest of the $P-3$ population members are derived by applying random mutations to these three base rescale factors, each with a mutation probability of $p$.

\textit{Monotonically non-decreasing constraint}. Once the initial population has been established, we evaluate LLM perplexity for every individual by applying the relevant RoPE rescale factors to the target LLM and then calculating the perplexity of the input $\mathbf{X}$. Individuals ranking in the top-k based on this metric are chosen as parents for the subsequent evolutionary steps. Due to the enormity of the search space, straightforward application of mutation and crossover operations may frequently result in suboptimal candidates, incurring excessive perplexity evaluations. This issue is exacerbated as $L'$ increases, since each perplexity assessment involves a computationally intensive inference process.

To tackle this issue, we enforce a constraint that the sampled RoPE rescaling factors must be monotonic and non-decreasing: $\lambda_i \le \lambda_{i+1}$. We only consider RoPE variants meeting this requirement for perplexity evaluation in the LLM, which notably narrows the search space and thus improves efficiency. Concretely, this stipulates that $\lambda_i$ progresses monotonically with respect to the RoPE dimension, where $i = 0, \ldots, 63$. This monotonicity with respect to dimension draws upon insights from NTK theory~\cite{jacot2018neural,tancik2020fourier,ntk}, where it is argued that dimensions associated with higher frequency (smaller $i$) benefit from less interpolation, resulting in smaller $\lambda_i$, whereas lower-frequency (higher $i$) dimensions can accommodate greater interpolation, corresponding to larger $\lambda_i$.

\begin{algorithm}[t]
	\small
	\caption{The search algorithm}
	\textbf{Input:} target LLM, input samples $\mathbf{X}$, population size $P$, mutation size $N_1$, crossover size $N_2$, maximum number of iterations $\mathcal{T}$, mutation probability $p$\\

	\begin{algorithmic}[1]
		\label{alg:search}
		\STATE Top-k $\leftarrow \phi$;	
		\STATE $\text{P}_0$ $\leftarrow$ \textit{Initialize\_population\_with\_optimization} ($P$, $\mathbf{X}$, $p$); 
		\FOR{$i$ from $1$ to $\mathcal{T}$}
		\STATE \textit{Compute\_perplexity} (LLM, $\text{P}_{i-1}$, $\mathbf{X}$);
		\STATE  Top-k $\leftarrow$ \textit{Update\_Topk} (Top-k, $\text{P}_{i-1}$);
		\STATE $\text{P}_{mutation}$ $\leftarrow$ \textit{Mutation\_with\_mono\_constraint} (Top-k, $N_1$, $p$); 
		\STATE $\text{P}_{crossover}$ $\leftarrow$ \textit{Crossover\_with\_mono\_constraint} (Top-k, $N_2$);
		\STATE $\text{P}_i$ $\leftarrow$ $\text{P}_{mutation}$ $\cup$ $\text{P}_{crossover}$ $\cup$ Top-k;
		\ENDFOR
		\STATE Output the element in Top-k with the minimum perplexity;
	\end{algorithmic}
\end{algorithm}

\textbf{8$\times$ extension without fine-tuning}. Through evolutionary search, our approach successfully discovers non-uniform RoPE rescale factors that maintain critical dimensions and positional information, thereby reducing the loss caused by interpolation. As illustrated in Fig.\ref{fig:non-ft}, this technique enables the context window of LLaMA2 to be scaled from 4k up to 32k without necessitating fine-tuning. In comparison, prior approaches like PI and non-uniform NTK or YaRN exhibit a significant increase in perplexity beyond a 2$\times$ context window extension.

\subsection{Extending LLM Context Window to 2048K}
\label{sec:extendto2048k}

\textbf{Progressive extension to 2048k}. In this section, we present our approach for expanding the context length of pre-trained LLMs from the standard 4k to beyond 2048k. Our non-uniform positional interpolation allows us to achieve an 8$\times$ extension without requiring any fine-tuning. However, when a much larger increase is needed (e.g., 512$\times$), fine-tuning becomes essential. A common strategy involves searching for suitable RoPE rescale factors for the desired 2048k context and then applying fine-tuning. Despite this, the process is hindered by the extremely high computational demands. Additionally, as our observations indicate, achieving stable fine-tuning outcomes is difficult when dealing with such large extension scales (refer to Appendix).

Fortunately, LongRoPE demonstrates efficacy when applied to both the baseline and further fine-tuned versions of extended LLMs. As such, we propose a resource-efficient, incremental approach capable of attaining the target length of 2048k through only 1k steps of fine-tuning, with the maximum training sequence restricted to 256k tokens.

$\Diamond$ \textit{Adapting pre-trained LLMs to 256k via LongRoPE search}. Using LLaMA2 as a representative case, we explore configurations aiming for context windows of 128k and 256k tokens. This process corresponds to extension ratios of 32$\times$ and 64$\times$, accordingly.

$\Diamond$ \textit{Fine-tuning to 256k}. Following pre-training, the LLM undergoes fine-tuning to reach a 256k context window size. Initially, we conduct 400 fine-tuning steps on LLaMA2 utilizing RoPE rescaled factors set for 128k. Subsequently, we switch the RoPE rescaled factors to 256k in the resulting checkpoint and perform a further 600 fine-tuning steps. This two-stage strategy demonstrates greater efficiency compared to fine-tuning directly with 256k.

$\Diamond$ \textit{Scaling fine-tuned extended LLM to 2048k via LongRoPE search}. In the final step, we conduct an additional search process on the already fine-tuned LLM with a 256k context length. Through this procedure, we are able to expand the context window size dramatically to 2048k, all without the need for additional fine-tuning. This leads to a total extension factor of $512\times$.

\textbf{Performance degradation within shorter context windows}. Upon expanding the context window to a considerable length of 2048k, we observe diminished model performance for inputs confined to the initial context window size. This phenomenon is a recognized limitation associated with positional interpolation~\cite{pi}, where the position embeddings corresponding to the original context must be compressed into a significantly reduced region in high-dimensional space, thereby impairing model efficacy. In scenarios where the context window is extended by a factor of 512, this crowding of position representations is especially pronounced for positions within the initial 4k range.

To address this issue, we conduct an additional evolutionary search on the extended LLM to fine-tune the RoPE rescale parameters specifically for shorter context lengths, such as 4k and 8k. For these shorter contexts, we restrict the upper bound of the searched $\lambda$, as they demand less positional interpolation. In the inference phase, the LLM adaptively modifies the RoPE rescale factors as needed.

 \section{Experiments}

\label{sec:exp}
\subsection{Setup}
\textbf{Evaluation Tasks and Models}. LongRoPE is incorporated into both LLaMA2-7B and Mistral-7B for experimental analysis. Our evaluation spans three primary areas: (1) measuring the perplexity of large language models with extended context on lengthy documents; (2) assessing the Passkey retrieval task, which tests the model's capacity to locate a straightforward passkey amid extensive distractor content; and (3) evaluating performance on widely-used LLM benchmarks under a standard 4096-token context window.

\textbf{Fine-tuning}. For LLaMA2, we adopt a learning rate of 2e-5 with a linear decay schedule and employ a global batch size of 32. The model is fine-tuned for 400 iterations on the Redpajama~\cite{redpajama} dataset, which is divided into 128k token segments, each demarcated by BOS and EOS tokens. Subsequently, building upon this checkpoint, training proceeds for an additional 600 steps to expand the context window to 256k tokens. Training with a 128k context length is conducted on 8 A100 GPUs utilizing the distributed framework~\cite{cube}; scaling up to 256k requires the use of 16 A100 GPUs. For Mistral, we use a fixed learning rate of 1e-6 and set the global batch size to 64. Both the 128k and 256k context models adopt the approach from YaRN~\cite{yarn}, training for 400 steps with 16k token sequences on Together Computer's Long-Data Collections~\cite{mistraldata}. This training leverages 4 A100 GPUs.

\textbf{Search}. When the target window size does not exceed 256k, the parameters are set to: $P$=64, $N_1$=$N_2$=16, $p$=0.3, and $\mathcal{T}$=40, while the top 32 candidates are chosen for mutation and crossover per iteration. For evaluation, perplexity is computed using five randomly selected samples from the PG19 validation dataset, ensuring that each sample meets the minimum required length equivalent to the target context. In cases where the window size surpasses 512k, both the population size and the number of mutations/crossovers are reduced by half. Here, perplexity assessment relies on three randomly selected instances from the Pile-Books3~\cite{pile} validation set.

\textbf{Baselines}. For achieving a 2048k context length, we performed fine-tuning on models with context windows of 128k and 256k tokens. This process results in LongRoPE-2048k (ft=128k) and LongRoPE-2048k (ft=256k), corresponding to the LLaMA2 and Mistral architectures. These four models are evaluated alongside leading context extension approaches, focusing on publicly available large language models that have undergone further fine-tuning following positional interpolation methods such as PI, NTK, and YaRN. The comparison set includes Together-32k~\cite{together}, Code LLaMA~\cite{codellama}, LongLoRA-full-FT-100k~\cite{longlora}, as well as YaRN-LLaMA and YaRN-Mistral~\cite{yarn}.

\subsection{Main Results}

\label{sec:mainresult}
\textbf{Language modeling for long sequences up to 256k.} Our analysis begins with a comparison against leading extended LLMs over sequences extending to 256k tokens. To illustrate the versatility of our approach, we conduct experiments using both the Proof-pile~\cite{pg19} and PG19~\cite{pile} test datasets. Perplexity is assessed across a range of context window sizes by employing a 256-token sliding window. For the PG19 dataset, the entire set of 100 test documents is included in our evaluation. For Proof-pile, mirroring the approach in YaRN~\cite{yarn}, we draw 10 samples at random, each sample containing at least 128k tokens.

Table~\ref{tbl:proof-pile} and Table~\ref{tbl:pg19} present a comparison of perplexity scores for LLaMA2 and Mistral models, each augmented with distinct interpolation strategies, evaluated on the Proof-pile and PG19 datasets. Two main findings emerge: \textbf{(1)} Our interpolated models demonstrate a consistent reduction in perplexity as the evaluation length increases from 4k to 256k, indicating an enhanced capacity to utilize extended contexts. \textbf{(2)} Notably, despite evaluating with context windows up to 16$\times$ longer than the default, which often hampers performance at shorter spans, our LongRoPE-2048k variants surpass leading baseline approaches at context lengths up to 256k.

\textbf{Long sequence language modeling beyond 2000k}. To assess performance on very lengthy documents, we utilize the Books3~\cite{pile} dataset. For efficient evaluation, we randomly choose 20 books, each longer than 2048k tokens, and apply a 256k sliding window approach.

As indicated in Table~\ref{tbl:books3}, LongRoPE effectively expands the context window of both LLaMA2-7B and Mistral-7B models up to 2048k, while maintaining perplexity values that match or outperform those of the baseline models at shorter sequence lengths ranging from 8k to 128k. A clear divergence in performance between the LLaMA2 and Mistral models is apparent at the 2048k context length. Notably, Mistral exhibits superior results compared to baselines for short contexts; however, its perplexity surpasses 7 for sequences longer than 256k. On the other hand, LLaMA2 demonstrates predictable performance behavior: its perplexity consistently improves with longer context windows, apart from minor increases at 1024k and 2048k. For LLaMA2, LongRoPE-2048k demonstrates optimal results when fine-tuned at 256k, which can be attributed to the reduced secondary extension ratio (8$\times$ as opposed to 16$\times$). Conversely, Mistral achieves better performance using a fine-tuning context window of 128k. This disparity is primarily due to the fact that, for both the 128k and 256k fine-tuning settings with Mistral, we utilize a 16k training length following the YaRN methodology, which consequently impacts Mistral’s capability to extend its context window post fine-tuning.

\textbf{Passkey retrieval}. Next, we examine the practical context window size for generative tasks. We adopt a synthetic benchmark for passkey retrieval introduced by~\cite{passkey}, in which the model must locate a randomly chosen five-digit number embedded somewhere within an extensive document. The exact prompt structure is provided in the appendix. For this experiment, we conduct 10 trials where, in each iteration, the passkey is positioned uniformly at random throughout the context length evaluated.

Fig.~\ref{fig:passkey} illustrates a comparative analysis of retrieval accuracy against baseline models. While previous approaches experience a steep decline in accuracy to zero once the input sequence exceeds 128k tokens, our LongRoPE-LLaMA2-2048k (ft=256k) consistently upholds high retrieval accuracy ($\ge$90\%) across the entire range from 4k to 2048k tokens, even when faced with the demanding task of locating a passkey in million-token contexts. Similarly, LongRoPE-Mistral-2048k (ft=128k) retains perfect (100\%) accuracy up to 1800k tokens before dropping to 60\% at 2048k, a trend that concurs with Table~\ref{tbl:books3}, where a modest increase in perplexity is noted at the 2048k mark.

\textbf{Conventional benchmarks applied to the original context window}. The performance of LongRoPE-2048k models is assessed within the initial context window through the Hugging Face Open LLM Leaderboard~\cite{huggingfaceboard}, considering both zero-shot and few-shot scenarios. The specific tasks include ARC-Challenge~\cite{allenai:arc} with 25 shots, HellaSwag~\cite{zellers2019hellaswag} with 10 shots, MMLU~\cite{mmlu} with 5 shots, and TruthfulQA~\cite{lin2021truthfulqa} in the zero-shot configuration.

As illustrated in Table~\ref{tbl:commonsense}, our models deliver similar performance to existing approaches on the original benchmark, which was developed for a more limited context window; notably, they exceed the original Mistral’s score on TruthfulQA by +0.5\%. While LongRoPE-LLaMA2-2048k, which was fine-tuned using a 256k context window, exhibits a marginally greater decline in performance, its results are still generally within acceptable limits across the majority of evaluated tasks.

\subsection{Ablation Results}

\textbf{Effectiveness of the second positional interpolation}. Our progressive extension approach applies a secondary, non-uniform positional interpolation to LLMs that have already been fine-tuned and extended, employing our search algorithm during this phase. To assess the impact of this step, we experiment with our fine-tuned LLaMA2-256k model and further expand its context window to 512k, 1024k, and 2048k tokens using both PI and YaRN methods. As reported in Table~\ref{tbl:secondsearch}, the use of our non-uniform positional interpolation maintains stable perplexity values even as the sequence length increases. In comparison, perplexity metrics with PI and YaRN rise sharply as the extension ratio grows.

\textbf{Recovery performance at short context lengths.} In order to address the decline in performance at reduced context sizes, we re-optimize the RoPE factors for LongRoPE-2048k by utilizing our search procedure. Concretely, we constrain the search to lower maximum scale factors, aiming to minimize interpolation for shorter contexts of 4k and 8k. 
Table~\ref{tbl:shortperformance} presents a perplexity assessment of LongRoPE-LLaMA2-2048k on Proof-pile at both 4k and 8k lengths, as well as average LLM benchmark scores. The data indicate that the approach yields notable gains at these shorter context windows.

\textbf{Analysis on the two forms of non-uniformities}. To disentangle the contributions of the two types of non-uniformities, we conduct ablation studies targeting each aspect individually. Our experimental design consists of two parts: (i) increasing the context window of LLaMA2-7B to 16k and 32k tokens using three strategies—PI, searching only for the optimal RoPE dimension, and jointly searching over both non-uniformities; (ii) further scaling our 256k-length fine-tuned LLaMA2 to a 2048k context, employing analogous methods. All evaluations of perplexity are performed without additional fine-tuning. From the results in Table~\ref{tbl:searchdimension}, we observe that introducing non-uniformity in the RoPE dimension considerably lowers perplexity relative to linear interpolation as used in PI. Modifying non-uniformity in token position yields performance benefits for 16k and 32k context lengths but does not exhibit notable advantages at the 2048k scale, potentially attributable to the challenges inherent in handling such extended sequences. Simply retaining the initial tokens without further interpolation proves ineffective at these lengths, which suggests an area for future investigation.

\section{Related Works}

In addition to methods based on position interpolation, this section discusses related works of other approaches. 

\textbf{Retrieval-based} methods incorporate an external memory component to store extended historical context, alongside retrieval mechanisms that fetch pertinent documents during inference~\cite{longllama,wang2023augmenting,borgeaud2022improving}. Such approaches often require substantial alterations to the LLM architectures. In comparison, our method remains more streamlined, necessitating only slight adjustments to positional embeddings. Moreover, our approach accommodates a broader range of long-context applications beyond simple retrieval, including long document summarization and few-shot learning scenarios.

\textbf{Attention-based context window extensions}. In addition to extending context windows via positional embedding interpolation, several studies extend the effective input length by adapting attention mechanisms while retaining the original context window size of LLMs~\cite{han2023lminfinite, streamingllm,parallelwindow}. Central to these approaches is the use of innovative attention masks designed to address the challenge of attention explosion arising from introducing novel token positions. Notably, such methods are compatible with, and enhance, positional interpolation techniques.

\textbf{Fine-tuning based approaches} aim to adapt pre-trained LLMs by altering position embeddings to accommodate extended sequence lengths. Techniques such as Code LLaMA~\cite{codellama}, LLaMA2 Long~\cite{xiong2023effective}, and ScaledRoPE~\cite{liu2023scaling} employ very large base values for RoPE and subsequently perform fine-tuning on the desired context length. In contrast, our method is versatile, supporting a range of target lengths and achieving context windows exceeding 2M tokens. Recently, as scaling to long contexts (e.g., above 128k) poses significant GPU memory challenges, approaches like LongLoRA~\cite{longlora} and PoSE~\cite{zhu2023pose} have been introduced to reduce the computational burden. Notably, our proposed method is complementary to these efficiency-driven fine-tuning strategies.

\section{Conclusion}

In this paper, we introduce LongRoPE, a technique that significantly increases the context window of LLMs to an unprecedented 2048k tokens, all while preserving their effectiveness within their initial, shorter context ranges. Our approach leverages two types of non-uniformities found in the RoPE positional embedding via a computationally efficient evolutionary search. This strategy yields two main advantages: it supplies robust initialization for subsequent fine-tuning and allows for an 8$\times$ expansion of the context window even in the absence of fine-tuning. Building upon these results, we develop a stepwise extension approach, whereby LLMs fine-tuned on 256k tokens are progressively adapted to handle a 2048k context window, eliminating the need for additional fine-tuning. Comprehensive empirical results demonstrate the robust performance of LongRoPE. We anticipate that the LongRoPE-2048k models will open new avenues for applications that require long context processing and foster further advancements in the field.

\section*{Broader Impacts} This research seeks to contribute to the development of the field of Machine Learning. While various societal impacts may arise from our findings, we do not identify any that require particular emphasis at this time.

	\balance

\end{document}